\chapter*{Lampiran}
\addcontentsline{toc}{chapter}{LAMPIRAN}

\section*{Lampiran 1. Gambaran Inovasi/TTG}
Skema arsitektur sistem Welya disajikan pada Gambar~\ref{gbr:arsitektur}.
Sistem terdiri atas antarmuka yang diakses pengguna melalui browser web
maupun aplikasi desktop (dibangun dengan Tauri), layanan API di sisi
server, komponen AI berbasis model bahasa besar, basis data, serta
integrasi resmi dengan layanan Google melalui otorisasi OAuth 2.0.

\begin{figure}[htbp]
  \centering
  \begin{tikzpicture}[
      font=\small,
      kotak/.style={draw=black!60, rounded corners, align=center,
      fill=blue!6, minimum height=1.3cm, text width=3.1cm},
      eksternal/.style={kotak, fill=green!8},
      data/.style={kotak, fill=orange!10},
      panah/.style={-{Stealth[length=2.5mm]}, semithick},
      ket/.style={font=\scriptsize, align=center, fill=white, inner sep=1pt}
    ]
    \node[kotak] (pengguna) {\textbf{Pengguna}\\browser web /\\aplikasi desktop};
    \node[kotak, right=1.1cm of pengguna] (frontend)
    {\textbf{Frontend}\\web React \emph{app.welya.biz.id}\\desktop (Tauri)};
    \node[kotak, right=1.1cm of frontend] (backend)
    {\textbf{Backend}\\API server Node.js\\\emph{api.welya.biz.id}};
    \node[eksternal, above=1.1cm of backend] (ai)
    {\textbf{Komponen AI}\\model bahasa besar\\(Azure OpenAI / Gemini)};
    \node[data, below=1.1cm of backend] (db)
    {\textbf{Basis Data}\\PostgreSQL};
    \node[eksternal, below=1.1cm of frontend] (google)
    {\textbf{Layanan Google}\\Gmail $\cdot$ Google Tasks\\Google Calendar};

    \draw[panah] (pengguna) -- node[ket, above] {HTTPS} (frontend);
    \draw[panah] (frontend) -- node[ket, above] {REST API} (backend);
    \draw[panah, <->] (backend) -- node[ket, right] {analisis email,\\subtask, estimasi} (ai);
    \draw[panah, <->] (backend) -- node[ket, right] {baca/tulis\\data} (db);
    \draw[panah, <->] (backend.south west) -- node[ket, sloped, above]
    {OAuth 2.0} (google.north east);
  \end{tikzpicture}
  \caption{Skema arsitektur sistem Welya}
  \label{gbr:arsitektur}
\end{figure}

Alur kerja sistem: pengguna menghubungkan akun Google melalui OAuth 2.0;
backend membaca email kampus dari Gmail secara berkala; komponen AI
menyaring email yang mengandung penugasan, mengekstraksi atribut tugas,
dan memecahnya menjadi subtask dengan estimasi waktu; hasilnya disimpan
di basis data, disinkronkan ke Google Tasks dan Google Calendar, lalu
ditampilkan pada dasbor pengguna.

Tampilan antarmuka Welya disajikan pada
Gambar~\ref{gbr:dasbor}--\ref{gbr:workspace}.

\begin{figure}[htbp]
  \centering
  \includegraphics[width=\textwidth]{dashboard-dark}
  \caption{Tampilan dasbor Welya: jadwal hari ini, tugas prioritas beserta
    progres dan subtask, serta rekomendasi dari AI berdasarkan waktu
  luang pengguna}
  \label{gbr:dasbor}
\end{figure}

\begin{figure}[htbp]
  \centering
  \includegraphics[width=\textwidth]{tasks-dark}
  \caption{Tampilan papan tugas (Tugas Saya) dengan status, prioritas,
  subtask, dan tenggat setiap tugas}
  \label{gbr:tugas}
\end{figure}

\begin{figure}[htbp]
  \centering
  \includegraphics[width=\textwidth]{calendar-dark}
  \caption{Tampilan kalender yang menyatukan jadwal kuliah dan tenggat
  tugas dalam satu tampilan}
  \label{gbr:kalender}
\end{figure}

\begin{figure}[htbp]
  \centering
  \includegraphics[width=\textwidth]{course-dark}
  \caption{Tampilan ruang mata kuliah: jadwal, tugas, materi, catatan,
  dan pengumuman setiap mata kuliah dalam satu tempat}
  \label{gbr:workspace}
\end{figure}

\clearpage
\section*{Lampiran 2. Perencanaan Teknis}

\subsection*{Rencana Anggaran Biaya (RAB)}
\begin{table}[htbp]
  \centering
  \caption{Rencana anggaran biaya}
  \small
  \begin{tabular}{@{}clrrr@{}}
    \toprule
    No & Uraian & Volume & Harga Satuan (Rp) & Jumlah (Rp) \\
    \midrule
    1 & Domain dan hosting (1 tahun)   & 1 &  &  \\
    2 & Kredit API AI                  & 1 &  &  \\
    3 & Kuota internet pengembangan    & 4 &  &  \\
    4 & Pendaftaran hak cipta (HKI)    & 1 &  &  \\
    5 & Cetak dan jilid laporan        & 3 &  &  \\
    \midrule
    \multicolumn{4}{r}{Total} &  \\
    \bottomrule
  \end{tabular}
\end{table}

\subsection*{Jadwal Kegiatan}
Pengerjaan dilaksanakan pada 3 Agustus--17 September 2026 dengan rincian
pada Tabel~\ref{tab:jadwal}.

\begin{table}[htbp]
  \centering
  \caption{Jadwal kegiatan}
  \label{tab:jadwal}
  \small
  \begin{tabular}{@{}clc@{}}
    \toprule
    No & Kegiatan & Waktu Pelaksanaan \\
    \midrule
    1 & Analisis kebutuhan                   & 3--8 Agustus 2026              \\
    2 & Perancangan sistem                   & 10--15 Agustus 2026            \\
    3 & Implementasi                         & 17 Agustus--5 September 2026   \\
    4 & Pengujian                            & 31 Agustus--11 September 2026  \\
    5 & Evaluasi \& penyusunan laporan       & 7--15 September 2026           \\
    6 & Artikel publikasi \& pendaftaran HKI & 14--17 September 2026          \\
    \bottomrule
  \end{tabular}
\end{table}

\clearpage
\section*{Lampiran 3. Dokumen Pendukung Mitra (khusus TTG)}
Surat pernyataan kerja sama, profil mitra, dan bukti kebutuhan mitra
(kondisi lapangan) dilampirkan di sini apabila tugas akhir dikategorikan
sebagai TTG dengan mitra.
